\documentclass[11pt,a4paper]{article}
\usepackage[utf8]{inputenc}
\usepackage[T1]{fontenc}
\usepackage[french]{babel}
\usepackage{amsmath,amssymb,amsfonts}
\usepackage{geometry}
\geometry{hmargin=2.5cm,vmargin=2.5cm}

\title{Exercice : Étude d'une fonction du second degré}
\author{}
\date{}

\begin{document}

\maketitle

Une entreprise fabrique et commercialise un objet connecté. Le coût total de production, exprimé en milliers d'euros, pour la fabrication de $x$ centaines d'objets (avec $x \in [0 ; 10]$) est modélisé par la fonction $C$ définie par :
\[C(x) = x^2 - 6x + 13\]

Chaque centaine d'objets est vendue au prix fixe de $4$ milliers d'euros. Le chiffre d'affaires total réalisé pour la vente de $x$ centaines d'objets est donc donné par la fonction $R(x) = 4x$.

\begin{enumerate}
    \item \textbf{Calculs préliminaires :}
    \begin{itemize}
        \item Quel est le coût total de production pour la fabrication de $2$ centaines d'objets ?
        \item Quel est le chiffre d'affaires réalisé pour la vente de $2$ centaines d'objets ? L'entreprise réalise-t-elle un bénéfice ou une perte pour cette quantité ?
    \end{itemize}

    \item \textbf{Étude du bénéfice :}
    \begin{itemize}
        \item On définit le bénéfice de l'entreprise (en milliers d'euros) par la fonction $B(x) = R(x) - C(x)$. Montrer que pour tout $x \in [0 ; 10]$, on a :
        \[B(x) = -x^2 + 10x - 13\]
    \end{itemize}

    \item \textbf{Forme canonique et extremum :}
    \begin{itemize}
        \item Mettre la fonction $B$ sous sa forme canonique.
        \item En déduire le bénéfice maximal que peut réaliser l'entreprise, ainsi que la quantité d'objets (en centaines) pour laquelle ce maximum est atteint.
    \end{itemize}

    \item \textbf{Seuil de rentabilité :}
    \begin{itemize}
        \item Déterminer algébriquement les quantités d'objets pour lesquelles l'entreprise réalise un bénéfice positif ou nul ($B(x) \ge 0$). Arrondir les valeurs au centième si nécessaire.
    \end{itemize}
\end{enumerate}

\end{document}